\subsection{某些归纳极限中的有界子集}

命题 5. --- 设 \((\mathrm{E}_{i})_{i\in\mathbf{I}}\) 是一个 Hausdorff 局部凸空间族，E 是这个族的拓扑直和（II, 第 29 页）。要使 E 的子集 B 有界，必须且只需存在 I 的一个有限子集 J，使得 \(\operatorname{pr}_{i}(\mathbf{B})\) 在 \(E_{i}\) 中有界（对 \(i\in J\) ），且 \(\operatorname{pr}_{i}(\mathbf{B})\subset\{0\}\) 对所有 \(i\notin J\) 成立。

设 J 是 I 的一个有限子集。由于 E 的拓扑在 \(\prod_{j\in J}E_{j}\) 上诱导乘积拓扑（II, 第 30 页, 命题 7 与第 31 页, 命题 8），由 III, 第 4 页, 推论 2 即知该条件是充分的。

反之，设 B 是 E 的一个有界子集。那么 \(\mathrm{pr}_{i}(\mathbf{B})\) 对所有 i 有界（III, 第 4 页, 推论 1）。因此只需证明存在 I 的一个有限子集 J，使得 \(\mathrm{pr}_{i}(\mathbf{B}) \subset \{0\}\) 对所有 \(i \notin J\) 成立。若不然，则存在 I 的互异元素的一个无穷序列 \((i_{n})\) 与 B 的元素的一个无穷序列 \((x_{n})\) ，使得 \(\mathrm{pr}_{i_{n}}(x_{n}) \neq 0\) 。由于 \(E_{i_{n}}\) 是 Hausdorff 的，存在一个连续半范数 \(p_{n}\) （\(E_{i_{n}}\) 上的），使得 \(p_{n}(\mathrm{pr}_{i_{n}}(x_{n})) \geqslant n\) 。于是 \(p = \sum_{n \geqslant 1} p_{n} \circ \mathrm{pr}_{i_{n}}\) 是 E 上的一个连续半范数，且 p 在 B 上不是有界的，这就得到矛盾。

命题 6. --- 设 E 是一个局部凸空间，它是 E 的闭向量子空间的一个递增序列 \((E_{n})\) 的严格归纳极限（II, 第 33 页）。E 的子集 B 有界，当且仅当它含于某个子空间 \(E_{n}\) 之中，且在这个子空间中有界。

该条件是充分的，因为 E 的拓扑在 \(E_{n}\) 上诱导的拓扑正是 \(E_{n}\) 上给定的拓扑（II, 第 32 页, 命题 9）。要看出该条件是必要的，只需（III, 第 4 页, 命题 3）证明：如果 E 的点的一个序列 \((x_{m})\) 不含于任何子空间 \(E_{n}\) 之中，那么它不能趋于 0。通过抽取序列 \((x_{m})\) 的一个子序列，我们可以假设存在一个严格递增的整数序列 \((n_{k})\) ，使得对每个指标 k，有 \(x_{k} \notin E_{n_{k}}\) 且 \(x_{k} \in E_{n_{k+1}}\) 。那么存在（II, 第 33 页, 引理 2）凸集的一个递增序列 \((V_{k})\) ，使得 \(V_{k}\) 是 \(E_{n_{k}}\) 中 0 的一个邻域，\(V_{k+1} \cap E_{n_{k}} = V_{k}\) ，且对每个指标 k 有 \(x_{k} \notin V_{k+1}\) 。于是诸 \(V_{k}\) 的并 V 是 E 中 0 的一个邻域，且对所有 k 有 \(x_{k} \notin V\) 。这就证明序列 \((x_{k})\) 不趋于 0。

对空间 \(\mathbf{E}\) （它是 \(\mathrm{E}\) 的闭子空间的一个不可数有向集的归纳极限），命题 6 的结论未必成立（参见 III, 第 38 页, 习题 7）。

命题 7. --- 设 \((\mathrm{E}_{n})_{n\geq0}\) 是一个 Hausdorff 局部凸空间序列，且对每个 n，设 \(u_{n}:E_{n}\to E_{n+1}\) 是一个单射的紧线性映射（即：存在 \(E_{n}\) 中 0 的一个邻域，它在 \(u_{n}\) 下的像是相对紧的；这蕴含 \(u_{n}\) 连续）。设 E 是系统 \((\mathrm{E}_{n},u_{n})\) 的归纳极限（II, 第 29 页），且设 \(v_{n}\) 是从 \(E_{n}\) 到 E 中的典范映射。那么局部凸空间 E 是 Hausdorff 的。此外，对 E 的每个子集 A，下列条件等价：

\begin{enumerate}
\def\labelenumi{(\roman{enumi})}
\item
  A 是有界的；
\item
  存在一个整数 \(n\) ，使得 \(\mathbf{A}\) 是在 \(v_{n}\) 之下 \(\mathrm{E}_n\) 的一个有界子集的像；
\item
  A 是相对紧的。
\end{enumerate}

我们把 \(\mathbf{E}_n\) 与 \(\mathbf{E}\) 的一个向量子空间视为同一（它赋有一个比诱导拓扑更细的拓扑）。

引理 1. --- 在命题 7 的假设下，E 的拓扑是使所有映射 \(v_{n} : \mathrm{E}_{n} \to \mathrm{E}\) 都连续的最细拓扑。

我们需要证明：如果 U 是 E 的一个子集，使得 \(U \cap E_{n}\) 对每个 n 在 \(E_{n}\) 中是开的，那么 U 在 E 中是开的；换言之，我们必须证明，对每个 \(x \in U\) ，存在一个均衡凸集 V，使得 \(x + V \subset U\) ，且 \(V \cap E_{n}\) 对所有充分大的 n 是 \(E_{n}\) 中 0 的一个邻域（II, 第 27 页, 命题 5）。对每个 n，设 \(W_{n}\) 是 \(E_{n}\) 中 0 的一个均衡凸邻域，使得 \(H_{n}\) （即 \(W_{n}\) 在 \(E_{n+1}\) 中的闭包）是紧的。设 \(x \in U\) ，并设 \(n_{0}\) 是一个使得 \(x \in E_{n_{0}}\) 的整数。我们将用归纳法构造标量的一个序列 \((\varepsilon_{n})_{n \geqslant 0}\) （\textgreater0），使得 \(x + \sum_{n_{0} \leqslant i \leqslant n} \varepsilon_{i}H_{i}\) 对 \(n \geqslant n_{0}\) 含于 U 中。假设诸 \(\varepsilon_{i}\) （对 i \textless{} n）已经构造出来。如果 \(n = n_{0}\) ，令 \(V_{n-1} = \{0\}\) ；否则令

\[
\mathrm{V} _ {n - 1} = \sum_ {n _ {0} \leqslant i \leqslant n - 1} \varepsilon_ {i} \mathrm{H} _ {i}.
\]

那么 \(V_{n-1}\) 在 \(E_{n}\) 中是紧的，从而在 \(E_{n+1}\) 中更是紧的。由于 \(U \cap E_{n+1}\) 在 \(E_{n+1}\) 中是开的，存在一个标量 \(\varepsilon_{n} > 0\) ，使得 \(x + V_{n} = x + V_{n-1} + \varepsilon_{n}H_{n}\) 含于 U 中（GT, II, §4, No.~3, cor.）。令 \(V = \bigcup_{n \geqslant n_{0}} V_{n}\) 。那么 V 是凸的且均衡的；对 \(n \geqslant n_{0}\) ，我们有 \(V \cap E_{n} \supset \varepsilon_{n}H_{n} \cap E_{n} \supset \varepsilon_{n}W_{n}\) ，因而 \(V \cap E_{n}\) 是 \(E_{n}\) 中 0 的一个邻域。这样就完成了引理的证明。

集合 \(\mathrm{U} = \mathrm{E} - \{0\}\) 具有这样的性质：集合 \(\mathrm{U} \cap \mathrm{E}_n = \mathrm{E}_n - \{0\}\) 在 \(\mathrm{E}_n\) 中是开的（对每个 \(n\) ），因而 \(\mathrm{U}\) 在 \(\mathrm{E}\) 中是开的，这就证明 \(\mathrm{E}\) 是 Hausdorff 的（GT, III, § 1, No.~3, prop. 2）。显然性质 (ii) 蕴含 (iii)，且 (iii) 蕴含 (i)。最后我们证明 (i) 蕴含 (ii)。为此只需证明：如果 \(\mathrm{A}\) （\(\mathrm{E}\) 的一个子集）不被诸集合 \(\sum_{0 \leqslant i \leqslant n} \mathrm{H}_i\) 中任何一个吸收，那么 \(\mathrm{A}\) 不是有界的。但这时存在一个序列 \((x_n)_{n \geqslant 1}\) （由 \(\mathrm{A}\) 的点组成），使得对每个 \(n\) ，有 \(x_n \notin n^2 \sum_{0 \leqslant i \leqslant n} \mathrm{H}_i\) 。

那么诸 \(x_{n}/n\) 所成之集依引理 1 是闭的，因为它与 \(E_{m}\) 的交对每个整数 m 都是离散的。于是诸 \(x_{n}/n\) 所成之集的补集是 0 的一个开邻域，它不吸收序列 \((x_{n})\) ，因而 A 不是有界的。

注. --- 1) 用上面的记号，设 \(F_{n}\) 是由 \(H_{n}\) 生成的向量空间，赋以等于 \(H_{n}\) 的度规的范数。我们将看到（III, 第 8 页, cor.），\(F_{n}\) 是一个 Banach 空间。从 \(F_{n}\) 到 \(E_{n+1}\) 中的单射是紧的，因而单射 \(w_{n}\) （从 \(F_{n}\) 到 \(F_{n+1}\) 中的单射）更是紧的。此外，E 是 Banach 空间的归纳系统 ( \(F_{n}, w_{n}\) ) 的归纳极限。事实上，E 中 0 的一个均衡凸邻域 V 满足：\(V \cap E_{n}\) 吸收 \(H_{n-1}\) （对所有 \(n \geqslant 1\) ）；反之，如果 E 中的一个均衡凸集 W 满足 \(W \cap E_{n}\) 吸收 \(H_{n-1}\) ，那么 \(W \cap E_{n-1}\) 含有一个与 \(W_{n-1}\) 位似的集合（对所有 \(n \geqslant 1\) ），从而 W 是 E 中 0 的一个邻域。

\begin{enumerate}
\def\labelenumi{\arabic{enumi})}
\setcounter{enumi}{1}
\tightlist
\item
  设 F 是一个局部凸空间，k 是一个整数 \(\geqslant0\) ，\(f:E^{k}\rightarrow F\) 是一个多重线性映射。要使 f 连续，必须且只需 f 在 \(E_{n}^{k}\) 上的限制对每个 n 连续。我们立即验证：\(E^{k}\) 对线性映射族 \(v_{n}\times\cdots\times v_{n}:E_{n}\times\cdots\times E_{n}\to E\times\cdots\times E\) 具有终局部凸拓扑（II, 第 28 页, cor. 2 与第 30 页, 命题 7），且 \(u_{n}\times\cdots\times u_{n}\) 是从 \((\mathrm{E}_{n})^{k}\) 到 \((\mathrm{E}_{n+1})^{k}\) 中的一个紧单射线性映射。现在只需应用引理 1。
\end{enumerate}
% semantic-fragments: legacy-input-seam
\relax{}
